\documentclass{beamer}

\usepackage[orientation=portrait, size=a0, scale=1.4]{beamerposter}
\usetheme{default} % Mejor para posters

\usepackage[spanish]{babel}
\usepackage[utf8]{inputenc}
\usepackage{graphicx} % Para imágenes
\usepackage{framed}   % Para cuadros
\usepackage{ragged2e} % Para mejor alineación

% Definición del comando \texto
\newcommand{\texto}[1]{\justifying #1}

\title{Título del Póster: Tema Principal de la Investigación}
\author{Autor1 A, Autor2 B, Autor3 C, Autor4 D}
\institute{Institute}
\date{Fecha}

\begin{document}
	
	\begin{frame}[t]
		
		\begin{columns}[t]
			
			% -------- Columna 1 --------
			\begin{column}{0.48\textwidth}
				
				\includegraphics[width=\textwidth]{fotos/Fig1poster.png}
				
				\begin{framed}
					Fig. 1. Descripción detallada de la figura. Incluir información relevante
					sobre lo que muestra la imagen, las condiciones experimentales y la escala.
				\end{framed}
				
			\end{column}
			
			% -------- Columna 2 --------
			\begin{column}{0.48\textwidth}
				
				\includegraphics[width=\textwidth]{fotos/Fig1poster.png}
				
				\begin{framed}
					Fig. 2. Descripción detallada de la figura.
				\end{framed}
				
			\end{column}
			
		\end{columns}
		
		\vspace{1cm}
		
		\texto{
			Praesent vitae porttitor tortor. Fusce mollis, lacus ut ullamcorper molestie,
			justo ligula congue velit, ut maximus magna lorem vitae arcu. Sed sodales justo
			nec faucibus ultricies. Nam dictum enim non venenatis egestas.
		}
		
	\end{frame}
	
\end{document}
